\documentclass[11pt]{article}

\usepackage{amsmath, amssymb, graphicx, geometry}
\geometry{margin=1in}

\title{Anisotropic Surface-Mapped Coarse-Grained Beads for Molecular Systems}
\author{VSEPR-SIM Project}
\date{\today}

\begin{document}

\maketitle

\begin{abstract}
Traditional coarse-grained (CG) models frequently reduce molecular fragments to isotropic point particles, preserving total mass and charge while discarding directional structure. This work proposes an intermediate representation: anisotropic surface-mapped beads, in which grouped atoms are projected to a single center while retaining directional interaction information via a surface descriptor. Using benzene as a canonical example, we demonstrate how this approach preserves planar symmetry and directional response while maintaining computational efficiency. The method provides a bridge between atomistic fidelity and reduced-order modeling suitable for multi-scale simulation.
\end{abstract}

\section{Introduction}

Coarse-grained modeling aims to reduce computational complexity by grouping atoms into larger interaction sites. Standard approaches represent each group as an isotropic particle with effective interaction potentials:

\begin{equation}
U_{ij} = U(r_{ij})
\end{equation}

While efficient, this approximation discards directional information inherent to many molecular structures. For planar or anisotropic systems such as benzene, this leads to loss of physically relevant behavior.

This work introduces a surface-mapped bead representation, preserving directional characteristics while retaining a single-site description.

\section{Atomistic to Coarse-Grained Mapping}

Let a bead $B$ be formed from atoms $i \in S_B$.

The center of mass is defined as:

\begin{equation}
\mathbf{R}_B = \frac{1}{M_B} \sum_{i \in S_B} m_i \mathbf{r}_i
\end{equation}

with total mass:

\begin{equation}
M_B = \sum_{i \in S_B} m_i
\end{equation}

and total charge:

\begin{equation}
Q_B = \sum_{i \in S_B} q_i
\end{equation}

A local coordinate frame is constructed using the inertia tensor of the atom group.

\section{Surface Descriptor Representation}

Instead of representing the bead as isotropic, we define a directional surface descriptor on a spherical shell:

\begin{equation}
S_B(\theta, \phi)
\end{equation}

where $(\theta, \phi)$ are angular coordinates in the bead's local frame.

This function encodes anisotropic properties such as:
\begin{itemize}
    \item effective density
    \item interaction strength
    \item electrostatic response
\end{itemize}

A practical formulation is obtained by sampling interaction with a probe particle:

\begin{equation}
U_B(r, \theta, \phi) = \sum_{i \in S_B} U_{i,p}(|\mathbf{x} - \mathbf{r}_i|)
\end{equation}

where $\mathbf{x} = \mathbf{R}_B + r \hat{n}(\theta, \phi)$.

\section{Compressed Representation}

To reduce storage cost, the surface descriptor may be expanded in spherical harmonics:

\begin{equation}
S_B(\theta, \phi) = \sum_{\ell,m} c_{\ell m} Y_{\ell m}(\theta, \phi)
\end{equation}

Low-order expansions are sufficient for many rigid molecular fragments.

\section{Application to Benzene}

Benzene is a planar molecule with strong anisotropy:

\begin{itemize}
    \item circular in-plane symmetry
    \item distinct out-of-plane behavior
    \item non-uniform interaction profile
\end{itemize}

A surface-mapped bead preserves:
\begin{itemize}
    \item planar geometry
    \item orientation-dependent interactions
    \item effective size and shape
\end{itemize}

while reducing the system from 12 atoms to a single interaction site.

\section{Interaction Model}

Interactions between beads are defined using directional dependence:

\begin{equation}
U = U(r, \Omega_B, \Omega_C)
\end{equation}

or via surface coupling:

\begin{equation}
U \approx \int S_B(\theta, \phi) S_C(\theta', \phi') K(r) \, d\Omega
\end{equation}

where $K(r)$ is a radial kernel.

\section{Dynamic Consistency and Updates}

The surface descriptor remains valid under small deformations. However, recomputation is required when:

\begin{itemize}
    \item bonds are broken or formed
    \item large conformational changes occur
    \item charge distribution changes significantly
\end{itemize}

\section{Advantages and Limitations}

\subsection{Advantages}
\begin{itemize}
    \item retains directional physics
    \item reduces atom count significantly
    \item enables efficient reuse of precomputed structure
\end{itemize}

\subsection{Limitations}
\begin{itemize}
    \item state dependence of effective interactions
    \item additional preprocessing cost
    \item need for descriptor compression
\end{itemize}

\section{Orientation-Coupled Interaction Framework}

The surface descriptor $S_B(\theta,\phi)$ encodes the intrinsic directional structure of an individual coarse-grained bead. However, this representation alone is insufficient to define interactions between beads, since intermolecular behavior depends not only on internal anisotropy but also on the \emph{relative orientation} between interacting units.

To construct a physically meaningful coarse-grained model, the interaction framework must therefore couple radial separation with orientational degrees of freedom, enabling the emergence of direction-dependent energetics such as stacking, alignment, and angular selectivity.

\subsection{Bead Orientation Representation}

Each bead $B$, formed from an underlying atomic subset, is endowed with a local orthonormal frame $Q_B = \{\hat{e}_1, \hat{e}_2, \hat{e}_3\}$ derived from the principal axes of its inertia tensor. This construction ensures that the orientation is both physically grounded and invariant under rigid-body transformations of the atomistic configuration.

For planar or near-planar structures, such as benzene, the orientation can be reduced to a primary unit vector $\hat{n}_B$, defined as the normal to the molecular plane. This provides a minimal yet sufficient representation of anisotropy for systems dominated by planar symmetry.

Given two interacting beads $A$ and $B$, define the relative displacement vector:

\begin{equation}
\mathbf{r} = \mathbf{R}_B - \mathbf{R}_A, \quad
r = |\mathbf{r}|, \quad
\hat{r} = \frac{\mathbf{r}}{r}
\end{equation}

The orientational relationship between the beads is captured through the scalar invariants:

\begin{equation}
\cos\theta_A = \hat{n}_A \cdot \hat{r}, \quad
\cos\theta_B = \hat{n}_B \cdot \hat{r}, \quad
\cos\phi = \hat{n}_A \cdot \hat{n}_B
\end{equation}

These quantities form a reduced set of rotationally invariant descriptors that fully characterize the relative alignment between two axisymmetric anisotropic beads.

\subsection{Reduced Interaction Model}

A tractable first-order anisotropic interaction potential may be constructed by augmenting a baseline radial interaction with orientation-dependent corrections:

\begin{equation}
U_{\text{eff}}(r,\theta_A,\theta_B,\phi)
=
U_0(r)
+
\lambda_1 f_1(r)\cos\phi
+
\lambda_2 f_2(r)\cos^2\theta_A
+
\lambda_3 f_3(r)\cos^2\theta_B
\end{equation}

Here, $U_0(r)$ represents an isotropic reference potential (e.g., Lennard-Jones or Morse), while the functions $f_i(r)$ modulate the strength of angular contributions as a function of separation. The coefficients $\lambda_i$ control the relative importance of different orientational modes.

This formulation introduces anisotropic sensitivity while preserving a low-dimensional parameter space, enabling efficient evaluation in large-scale simulations. Importantly, the chosen angular terms correspond to physically interpretable alignment tendencies, such as parallel orientation ($\cos\phi$) and alignment with the intermolecular axis ($\cos\theta$ terms).

\subsection{Surface Descriptor Coupling}

While the reduced model provides computational efficiency, higher-fidelity interactions may be derived directly from the surface descriptors themselves. In this formulation, the interaction energy arises from a convolution-like coupling between directional surface fields:

\begin{equation}
U \approx \int S_A(\theta,\phi)\, S_B(\theta',\phi')\, K(r)\, d\Omega
\end{equation}

where $K(r)$ is a radial kernel governing distance dependence, and the angular coordinates $(\theta',\phi')$ are obtained by transforming bead $B$'s local frame into that of bead $A$.

This representation effectively treats each bead as a distributed interaction surface rather than a point object, allowing fine angular structure---including multipolar features and localized interaction regions---to contribute directly to the interaction energy.

In practice, this integral may be evaluated using truncated spherical harmonic expansions, reducing the interaction to a finite set of mode couplings.

\subsection{Torque and Rotational Dynamics}

Because the interaction potential depends explicitly on orientation, translational forces alone do not fully describe system dynamics. Rotational degrees of freedom must be incorporated to ensure physical consistency.

The torque acting on bead $B$ is given by:

\begin{equation}
\boldsymbol{\tau}_B = -\frac{\partial U}{\partial \Omega_B}
\end{equation}

where $\Omega_B$ denotes the rotational configuration of the bead, parameterized for example by Euler angles or quaternions.

This torque induces rotational motion according to rigid-body dynamics, enabling phenomena such as alignment, libration, and orientational relaxation to emerge naturally from the interaction model.

\subsection{Application to Benzene Dimer Configurations}

The orientation-dependent framework is particularly relevant for systems such as benzene dimers, where interaction energy is strongly modulated by relative alignment.

The model captures several well-known configurations:

\begin{itemize}
    \item \textbf{Face-to-face stacking:} $\cos\phi \approx 1$, $\theta_A \approx \theta_B \approx 0$
    \item \textbf{Edge-to-face (T-shaped):} $\theta_A \approx \pi/2$, $\theta_B \approx 0$
    \item \textbf{Offset stacking:} intermediate angular values with lateral displacement
\end{itemize}

These configurations correspond to distinct energetic minima arising from $\pi$-$\pi$ interactions and electrostatic anisotropy. Such behavior cannot be reproduced within isotropic coarse-grained models, highlighting the necessity of orientation-aware formulations.

\subsection{Model Hierarchy}

The proposed framework supports a systematic hierarchy of interaction models with increasing fidelity:

\begin{itemize}
    \item \textbf{Isotropic:} $U(r)$
    \item \textbf{Axisymmetric anisotropic:} $U(r,\theta_A,\theta_B,\phi)$
    \item \textbf{Descriptor-resolved:} $U(r,\Omega_A,\Omega_B,S_A,S_B)$
\end{itemize}

This hierarchy enables adaptive resolution strategies, where simpler models are used for large-scale exploration and more detailed representations are invoked when directional effects become dominant.

In this sense, the orientation-coupled interaction framework serves as a bridge between classical coarse-grained potentials and fully resolved anisotropic interaction fields.

\subsection{Descriptor Enrichment for Complex Anisotropic Structures}

The spherical harmonic representation

\begin{equation}
S_B(\theta,\phi) = \sum_{\ell,m} c_{\ell m} Y_{\ell m}(\theta,\phi)
\end{equation}

provides a compact and efficient description of moderately anisotropic systems. However, for highly structured and sterically complex entities---such as coordination clusters with bulky ligands---low-order expansions are insufficient. In these cases, the angular variation of interaction properties is sharp, non-uniform, and often dominated by localized features.

\subsubsection{Higher-Order Angular Resolution}

To accurately capture directional structure, the spherical harmonic expansion must be extended to higher orders:

\begin{equation}
S_B(\theta,\phi) = \sum_{\ell=0}^{\ell_{\max}} \sum_{m=-\ell}^{\ell} c_{\ell m} Y_{\ell m}(\theta,\phi)
\end{equation}

where $\ell_{\max}$ is chosen based on the angular complexity of the system. Larger values of $\ell_{\max}$ allow representation of sharp steric occlusion and narrow interaction channels, at the cost of increased computational overhead.

\subsubsection{Multi-Channel Surface Representation}

A single scalar descriptor is generally insufficient to encode the diverse physical interactions present in complex systems. Instead, the surface descriptor is extended to a vector-valued representation:

\begin{equation}
S_B(\theta,\phi) \rightarrow
\left\{
S_B^{(\text{steric})}(\theta,\phi),\,
S_B^{(\text{elec})}(\theta,\phi),\,
S_B^{(\text{disp})}(\theta,\phi)
\right\}
\end{equation}

Each channel captures a distinct physical contribution:

\begin{itemize}
    \item $S^{(\text{steric})}$: geometric accessibility and excluded volume
    \item $S^{(\text{elec})}$: electrostatic potential or charge distribution
    \item $S^{(\text{disp})}$: dispersion and non-bonded interaction strength
\end{itemize}

This decomposition enables independent modeling and weighting of interaction mechanisms, improving both interpretability and accuracy.

\subsubsection{Full Orientation Representation}

For systems with three-dimensional anisotropy, a single orientation vector is insufficient. Instead, each bead is associated with a full orthonormal frame:

\begin{equation}
Q_B = \{\hat{e}_1, \hat{e}_2, \hat{e}_3\}
\end{equation}

derived from the principal axes of the inertia tensor. This allows the surface descriptor to be consistently defined in the bead's local coordinate system and properly transformed during interaction evaluation.

\subsubsection{Implications for Interaction Modeling}

The enriched descriptor directly affects the interaction formulation. The effective interaction between two beads becomes:

\begin{equation}
U \approx \int
\sum_k S_A^{(k)}(\theta,\phi)\, S_B^{(k)}(\theta',\phi')\, K_k(r)\, d\Omega
\end{equation}

where $k$ indexes interaction channels and $K_k(r)$ are channel-specific radial kernels.

This formulation enables:
\begin{itemize}
    \item directional steric screening
    \item anisotropic electrostatic interactions
    \item orientation-dependent dispersion
\end{itemize}

\subsubsection{Modeling Considerations}

The introduction of higher-order and multi-channel descriptors increases both representational power and computational cost. Practical implementations must therefore balance:

\begin{itemize}
    \item angular resolution ($\ell_{\max}$)
    \item number of descriptor channels
    \item evaluation efficiency
\end{itemize}

Adaptive schemes, in which descriptor complexity is increased only for highly anisotropic systems, may provide an effective compromise.

\subsubsection{Physical Interpretation}

At this level of representation, the bead is no longer a point particle with an effective potential, but rather a compressed representation of a three-dimensional interaction field. The modeling objective is therefore not simply reduction, but controlled preservation of physically relevant directional structure.

\section{Model Selection Guidelines}

To balance computational efficiency and physical fidelity, the interaction model is applied using a two-tier hierarchical strategy based on structural regime rather than a fixed atom-count threshold.

\subsection{Tier 1: Reduced Anisotropic Model}

The reduced orientation-coupled model is used for systems with low angular complexity. This includes structures that are:

\begin{itemize}
    \item rigid and compact
    \item smoothly anisotropic
    \item well-described by a single dominant axis or simple frame
    \item not dominated by localized interaction sites
\end{itemize}

Typical examples include aromatic ring systems such as benzene, naphthalene, pyridine, and phenol.

The interaction is represented using a low-dimensional orientation-dependent potential:

\begin{equation}
U_{\text{eff}}(r,\theta_A,\theta_B,\phi)
\end{equation}

This model is appropriate when:
\begin{itemize}
    \item low-order spherical harmonic expansion is sufficient
    \item a small number of angular invariants captures the interaction behavior
    \item computational efficiency is prioritized
\end{itemize}

\subsection{Tier 2: Enriched Descriptor Model}

The enriched descriptor-based model is used for systems with high angular complexity and heterogeneous interaction surfaces. This includes structures that are:

\begin{itemize}
    \item strongly three-dimensional
    \item sterically shielded or directionally gated
    \item chemically heterogeneous
    \item metal-centered or coordination-driven
    \item characterized by non-uniform surface accessibility
\end{itemize}

Typical examples include organometallic clusters, coordination complexes, and bulky substituted systems.

The bead is represented using a multi-channel surface descriptor:

\begin{equation}
S_B(\theta,\phi) \rightarrow
\left\{
S^{(\text{steric})},\,
S^{(\text{elec})},\,
S^{(\text{disp})}
\right\}
\end{equation}

with a full orientation frame:

\begin{equation}
Q_B = \{\hat{e}_1, \hat{e}_2, \hat{e}_3\}
\end{equation}

This model is required when:
\begin{itemize}
    \item interaction depends on angular accessibility
    \item low-order angular representations produce large error
    \item multiple physical interaction modes must be resolved
\end{itemize}

\subsection{Selection Heuristic}

As a practical guideline, the enriched model is typically activated for systems with:
\begin{itemize}
    \item atom count $N \gtrsim 18$
    \item presence of a metal center
\end{itemize}

However, the definitive criterion is the angular complexity of the surface descriptor. Systems that are poorly approximated by low-order angular representations must be promoted to the enriched model regardless of size or composition.

\subsection{Summary}

The reduced model provides efficiency for structurally simple systems, while the enriched model preserves directional physics in complex environments. This hierarchical approach enables adaptive resolution across a wide range of molecular and materials systems.

\section{Unified Descriptor Strategy}

The preceding model selection guidelines introduced a two-tier hierarchical strategy, distinguishing between a reduced orientation-coupled model and an enriched multi-channel descriptor model. While conceptually useful, this formulation relies on a discrete architectural switch---triggered by atom count, metal center presence, or angular complexity---that introduces a discontinuous boundary between model regimes.

This section presents a refinement: rather than maintaining two separate data models with distinct interaction machinery, all systems are represented within a \emph{single descriptor-based framework}. Model complexity is adjusted through controlled truncation rather than by changing the underlying formalism. The result is a continuous, unified architecture in which the previous ``tiers'' emerge as natural truncation levels of one model family.

\subsection{Motivation}

Hard switching between separate coarse-grained interaction models introduces several undesirable properties:

\begin{itemize}
    \item \textbf{Discontinuity:} The energy surface is non-smooth at the model boundary, complicating optimization and dynamics.
    \item \textbf{Non-transferability:} Parameters calibrated in one model regime do not apply in the other.
    \item \textbf{Error opacity:} It is difficult to compare model quality across regimes when the representations differ.
    \item \textbf{Arbitrary thresholds:} Heuristic cutoffs (e.g., $N \geq 18$, metal center present) do not correspond to a smooth physical variation.
\end{itemize}

A unified formalism avoids these issues by ensuring that model enrichment occurs within a single continuous framework.

\subsection{Unified Formalism}

All beads, regardless of structural complexity, are represented using the same multi-channel spherical harmonic descriptor:

\begin{equation}
S_B^{(k)}(\theta,\phi) = \sum_{\ell=0}^{\ell_{\max}} \sum_{m=-\ell}^{\ell} c_{\ell m}^{(k)} Y_{\ell m}(\theta,\phi)
\end{equation}

where $k$ indexes the interaction channel (steric, electrostatic, dispersion) and $\ell_{\max}$ controls angular resolution.

Each bead carries a full orientation frame $Q_B = \{\hat{e}_1, \hat{e}_2, \hat{e}_3\}$ and a set of descriptor coefficients $\{c_{\ell m}^{(k)}\}$ for each active channel.

The interaction energy between two beads is defined by a single universal form:

\begin{equation}
U = F\!\left(r,\, Q_A,\, Q_B,\, \{c_{\ell m}^{(k)}\}\right)
\end{equation}

which evaluates as:

\begin{equation}
U = U_0(r) + \sum_k \lambda_k \left(\sum_{\ell,m} c_{\ell m}^{(k)}(A)\, c_{\ell m}^{(k)}(B)\right) K_k(r)
\end{equation}

where $U_0(r)$ is the isotropic baseline, $K_k(r)$ are channel-specific radial kernels, and $\lambda_k$ are coupling weights.

This formulation subsumes both the reduced and enriched models as special cases:

\begin{itemize}
    \item \textbf{Reduced regime:} small $\ell_{\max}$, single active channel, few nonzero coefficients
    \item \textbf{Enriched regime:} larger $\ell_{\max}$, multiple active channels, richer coefficient structure
\end{itemize}

The transition between these regimes is not a model switch, but a change in the number of active coefficients within the same mathematical object.

\subsection{Adaptive Truncation}

Rather than selecting a model tier via fixed heuristic thresholds, the unified framework determines the appropriate resolution level through \emph{descriptor residual analysis}.

Given a bead's atomistic structure, the procedure is:

\begin{enumerate}
    \item Fit a low-order descriptor ($\ell_{\max} = \ell_0$)
    \item Compute the reconstruction residual $\varepsilon(\ell_0)$
    \item If $\varepsilon > \varepsilon_{\text{tol}}$, increase $\ell_{\max}$ and activate additional channels
    \item Repeat until residual is below tolerance or maximum resolution is reached
\end{enumerate}

The reconstruction residual is defined as:

\begin{equation}
\varepsilon(\ell_{\max}) = \frac{\sum_{i} \left| S_{\text{ref}}(\hat{n}_i) - S_{\ell_{\max}}(\hat{n}_i) \right|^2}{\sum_{i} \left| S_{\text{ref}}(\hat{n}_i) \right|^2}
\end{equation}

where $S_{\text{ref}}$ is the probe-sampled reference surface and $S_{\ell_{\max}}$ is the truncated spherical harmonic reconstruction.

This approach ensures that complexity is added continuously in response to representational need, rather than through fixed structural cutoffs.

\subsection{Resolution Levels}

The unified descriptor supports a continuous hierarchy of resolution levels, which replaces the discrete tier concept:

\begin{itemize}
    \item \textbf{Level 0 (Isotropic):} Only $c_{00}$ retained; equivalent to a point particle with effective radius
    \item \textbf{Level 1 (Axial):} $\ell_{\max} \leq 2$, single channel; captures dominant planar or axial anisotropy
    \item \textbf{Level 2 (Moderate):} $\ell_{\max} \leq 4$, one or two channels; sufficient for rigid aromatic systems
    \item \textbf{Level 3 (Enriched):} $\ell_{\max} \leq 8$, all channels active; required for sterically complex or metal-centered systems
\end{itemize}

This hierarchy is not a set of discrete models but a continuum of truncation depths within one mathematical framework.

\subsection{Advantages over Discrete Tier Selection}

The unified descriptor strategy provides several advantages:

\begin{itemize}
    \item \textbf{Smooth energy surface:} No discontinuities at model boundaries
    \item \textbf{Single optimization space:} Parameters, coefficients, and errors are all defined in the same representation
    \item \textbf{Measurable error:} Reconstruction residual provides a quantitative basis for resolution selection
    \item \textbf{Hierarchical refinement:} Start coarse, refine only where needed
    \item \textbf{Transferability:} The same descriptor object and energy machinery apply across all chemical systems
\end{itemize}

\subsection{Relationship to Previous Tier Concept}

The two-tier model selection guidelines presented in the preceding section remain valid as a practical engineering heuristic. However, they are now understood as approximate indicators of where a system falls within the continuous resolution hierarchy, rather than as triggers for architectural switching.

In implementation, the ``reduced'' regime corresponds to a low-order truncation of the unified descriptor, and the ``enriched'' regime corresponds to a higher-order truncation with additional active channels. The underlying data model and energy machinery are identical in both cases.

\section{Atomistic Preparation Layer}

The unified descriptor strategy establishes a mathematically coherent coarse-grained formalism. However, the quality of any coarse-grained representation is bounded by the quality of the atomistic input from which it is derived. If the atomistic preparation is unreliable---with incomplete topology, missing charges, or ill-defined local frames---the descriptor machinery inherits this uncertainty and the resulting coarse-grained model is untrustworthy regardless of its mathematical sophistication.

This section formalizes the atomistic preparation layer as a necessary prerequisite for reliable descriptor generation.

\subsection{Purpose}

The atomistic namespace exists to answer a fundamentally different question than the coarse-grained namespace:

\begin{itemize}
    \item \textbf{Atomistic:} ``What is this structure, exactly?''
    \item \textbf{Coarse-grained:} ``How do I compress this structure into a transferable reduced object?''
\end{itemize}

The atomistic layer is the source-of-truth layer that owns:

\begin{itemize}
    \item atoms, bonds, and formal charges
    \item local geometry and coordination environment
    \item ring detection and aromaticity flags
    \item molecular graph and canonical fragment decomposition
    \item local frames derived from real structure
\end{itemize}

Its outputs are consumed by the coarse-grained layer, but the reverse dependency must never exist. The atomistic layer must not depend on or import any coarse-grained types.

\subsection{Scale Boundary: The Fragment View}

Rather than allowing the coarse-grained layer to rummage through raw atom lists, a clean bridge object is defined at the scale boundary. This object---the \emph{fragment view}---packages all atomistic information needed for descriptor construction into a single, validated interface.

The fragment view contains:

\begin{itemize}
    \item atom records: position, atomic number, charge, flags
    \item bond records: atom indices and bond order
    \item aggregate properties: total mass, total charge
    \item a local coordinate frame derived from the inertia tensor
    \item structural classification flags: aromatic, cyclic, organometallic
\end{itemize}

The coarse-grained layer receives a fragment view and constructs a bead and its unified descriptor from it. This enforces a one-directional data flow:

\begin{equation}
\text{atomistic structure} \;\xrightarrow{\text{fragment view}}\; \text{unified descriptor} \;\rightarrow\; \text{bead interactions}
\end{equation}

\subsection{Validation and Status Codes}

The atomistic preparation layer must produce either a valid fragment view or an explicit failure with a reason code. Silent propagation of invalid inputs is prohibited.

Failure modes include:

\begin{itemize}
    \item invalid geometry (degenerate or collapsed atom positions)
    \item missing charges (incomplete charge assignment)
    \item ill-defined frame (insufficient atom count or degeneracy)
    \item unsupported topology (unrecognized bonding pattern)
\end{itemize}

Each failure mode produces a specific status code, making debugging tractable and preventing the coarse-grained layer from inheriting garbage.

\subsection{Architectural Implications}

The introduction of a formal scale boundary has several consequences:

\begin{itemize}
    \item \textbf{Dependency asymmetry:} The atomistic namespace never imports from \texttt{coarse\_grain}. The coarse-grained namespace depends on atomistic bridge types.
    \item \textbf{Testability:} The fragment view can be independently validated before descriptor construction.
    \item \textbf{Reusability:} The atomistic preparation layer supports additional consumers beyond coarse-graining: reference calculations, backmapping, validation against exact structure, fragment classification, and hybrid atomistic/CG regions.
    \item \textbf{Error traceability:} If a descriptor is incorrect, the fragment view provides the exact atomistic input that produced it.
\end{itemize}

\subsection{Relationship to Descriptor Generation}

The unified descriptor is constructed from a fragment view through the following chain:

\begin{enumerate}
    \item Atomistic layer produces a validated \texttt{FragmentView}
    \item The fragment view provides positions, charges, and frame to the mapping layer
    \item The mapping layer generates probe-based surface fields
    \item Surface fields are projected onto spherical harmonics to produce descriptor coefficients
    \item The reconstruction residual determines whether the current truncation is sufficient
\end{enumerate}

This chain ensures that every descriptor coefficient is traceable back through the fragment view to specific atoms in the original structure. The anti-black-box contract is maintained across the scale boundary.

\section{Anisotropic Bead Model --- Implementation Specification}

This section formalises the unified coarse-grained descriptor framework with
adaptive spherical harmonic resolution.  All bead instances conform to a single,
rigidly defined state structure; no dual-model branching is permitted.

\subsection{Core Data Structure}

A bead $B$ is defined by the canonical tuple
\begin{equation}
    B = \bigl\{ \mathbf{R}_B,\; \mathbf{Q}_B,\; \{c_{\ell m}^{(k)}\},\; M_B,\; Q_B^{\mathrm{tot}} \bigr\}
\end{equation}
where $\mathbf{R}_B$ is the centre of mass, $\mathbf{Q}_B$ is an orthonormal
$3\times3$ orientation frame, $\{c_{\ell m}^{(k)}\}$ are per-channel real
spherical harmonic coefficients, $M_B$ is the total mass, and
$Q_B^{\mathrm{tot}}$ is the total charge.

Three physical channels are defined:
\begin{center}
\begin{tabular}{cl}
    $k=0$ & Steric (excluded-volume / atomic density) \\
    $k=1$ & Electrostatic (Coulomb potential at surface) \\
    $k=2$ & Dispersion (Lennard-Jones field at surface) \\
\end{tabular}
\end{center}

\noindent
Invariants:
\begin{itemize}
    \item $\mathbf{Q}_B$ must remain orthonormal; enforce via Gram--Schmidt at every mutation site.
    \item $\ell_{\max}$ must be consistent per channel within a given bead.
    \item Memory layout must be contiguous to enable SIMD vectorisation in downstream kernels.
\end{itemize}

\subsection{Descriptor Generation Pipeline}

The pipeline maps a fully specified atomistic fragment to a fully populated
bead.  All steps are deterministic; stochastic initialisation is not permitted.

\begin{enumerate}
    \item \textbf{Centre computation:} $\mathbf{R}_B = M_B^{-1} \sum_{i \in S_B} m_i \mathbf{r}_i$
    \item \textbf{Local frame:} $\mathbf{Q}_B$ via inertia tensor diagonalisation (Jacobi)
    \item \textbf{Channel projection:} For each channel $k$, sample the surface field and project onto real spherical harmonics: $c_{\ell m}^{(k)} = \int S^{(k)}(\hat{n})\, Y_{\ell m}(\hat{n})\, d\Omega$
    \item \textbf{Scalar aggregation:} $M_B = \sum m_i$, \; $Q_B^{\mathrm{tot}} = \sum q_i$
\end{enumerate}

\subsection{Channel Sampling Definitions}

\begin{center}
\begin{tabular}{ll}
    \textbf{Channel} & \textbf{Physical Definition} \\
    \hline
    Steric        & Excluded-volume / atomic density projected onto bead surface \\
    Electrostatic & Coulomb potential evaluated at surface points from partial charges \\
    Dispersion    & Lennard-Jones field contribution at bead surface \\
\end{tabular}
\end{center}

\subsection{Interaction Engine}

The pairwise interaction energy between beads $A$ and $B$ is evaluated in
harmonic space.  Runtime surface integration is not employed.

\begin{equation}
    U_{AB}(r, \mathbf{Q}_A, \mathbf{Q}_B) = \sum_k I_{AB}^{(k)}(r, \mathbf{Q}_A, \mathbf{Q}_B)
\end{equation}

The per-channel interaction requires rotating the coefficients of $B$ into
the local frame of $A$ via Wigner-$D$ matrices for real spherical harmonics,
then summing:

\begin{equation}
    I_{AB}^{(k)} = \sum_{\ell=0}^{\ell_{\max}} \sum_{m=-\ell}^{\ell}
        c_{\ell m}^{(k),A}\; \tilde{c}_{\ell m}^{(k),B}\; K_k(\ell, r)
\end{equation}

where $\tilde{c}_{\ell m}^{(k),B}$ denotes the rotated coefficients and
$K_k(\ell, r)$ is the per-channel, per-$\ell$ radial kernel:

\begin{equation}
    K_{\mathrm{steric}}(\ell, r) = \frac{e^{-\alpha_s r}}{1 + \ell}
\end{equation}
\begin{equation}
    K_{\mathrm{electrostatic}}(\ell, r) = \frac{1}{r^{\ell+1}}
\end{equation}
\begin{equation}
    K_{\mathrm{dispersion}}(\ell, r) = \frac{-C_6}{r^{6+\ell}}
\end{equation}

\subsection{Adaptive Complexity}

Resolution is adjusted per-bead per-channel by measuring reconstruction
fidelity and incrementing $\ell_{\max}$ when the error metric exceeds a
threshold.  This eliminates dual-model architectures; a single model instance
spans the full complexity range.

\begin{equation}
    \varepsilon_B^{(k)} = \frac{\sum_i \bigl| S_{\mathrm{ref}}^{(k)}(\hat{n}_i) - \tilde{S}^{(k)}(\hat{n}_i) \bigr|^2}
                               {\sum_i \bigl| S_{\mathrm{ref}}^{(k)}(\hat{n}_i) \bigr|^2}
\end{equation}

If $\varepsilon_B^{(k)} > \varepsilon_{\mathrm{tol}}$, then $\ell_{\max}^{(k)} \leftarrow \ell_{\max}^{(k)} + \Delta\ell$ and the channel is reprojected.

Expected behaviour:
\begin{itemize}
    \item Small or symmetric molecules converge at low $\ell_{\max}$.
    \item Large or anisotropic structures self-select higher harmonic resolution.
    \item No branching into separate model classes---only resolution scaling.
\end{itemize}

\subsection{System-Level Integration}

The components specified above constitute a self-consistent physical model.
Compatible integration targets include:
\begin{itemize}
    \item MD integrators (velocity Verlet, Langevin)
    \item Monte Carlo sampling (Metropolis--Hastings)
    \item Bulk system simulations at arbitrary $N$
\end{itemize}

Required next steps:
\begin{enumerate}
    \item Precompute Wigner-$D$ rotation operators; cache per angular momentum order
    \item Implement neighbour lists to reduce interaction evaluation from $O(N^2)$ to $O(N)$
    \item Validate against benzene $\pi$-stacking reference energies
    \item Validate against naphthalene crystal packing geometry
    \item Validate against organometallic coordination geometries
\end{enumerate}

\section{Formal Bead State Decomposition}

The canonical bead state
\begin{equation}
    B = \bigl\{ \mathbf{R}_B,\; \mathbf{Q}_B,\; \{c_{\ell m}^{(k)}\},\; M_B,\; Q_B^{\mathrm{tot}} \bigr\}
\end{equation}
is not merely a tuple of numbers.  Each variable occupies a distinct physical
role, obeys a distinct update rule, and serves a distinct architectural
purpose.  This section decomposes the state into its constituent layers
and establishes the ownership, mutability, and numerical responsibilities
of every term.

\subsection{$\mathbf{R}_B$: Translational Anchor}

The centre-of-mass position $\mathbf{R}_B \in \mathbb{R}^3$ is the bead's
spatial location in the global (lab) frame.  It is computed at mapping time as
\begin{equation}
    \mathbf{R}_B = \frac{1}{M_B} \sum_{i \in S_B} m_i \,\mathbf{r}_i
\end{equation}
and thereafter evolves under the equations of motion.

\paragraph{Physical role.}
$\mathbf{R}_B$ is the translational anchor from which every pairwise
interaction begins.  It enters:
\begin{itemize}
    \item separation vectors: $\mathbf{r}_{AB} = \mathbf{R}_B - \mathbf{R}_A$,
          $r = |\mathbf{r}_{AB}|$
    \item neighbour-list construction and periodic image selection
    \item translational kinetic energy:
          $T_{\mathrm{trans}} = \tfrac{1}{2} M_B \dot{\mathbf{R}}_B^2$
    \item the atomistic$\to$CG mapping as the image of the fragment centroid
\end{itemize}

\paragraph{Update frequency.}
$\mathbf{R}_B$ is the most rapidly changing state variable.  It is updated
\emph{every dynamics step} by the integrator.

\paragraph{Ownership.}
$\mathbf{R}_B$ belongs to the \emph{kinematic shell} of the bead---the
layer concerned with motion through space, not with intrinsic identity.

\subsection{$\mathbf{Q}_B$: Orientation Frame}

The orientation $\mathbf{Q}_B \in \mathrm{SO}(3)$ is a $3\times 3$
orthonormal matrix whose columns are the bead's principal axes, ordered by
ascending eigenvalue of the inertia tensor.  It satisfies the invariant
\begin{equation}
    \mathbf{Q}_B^T \mathbf{Q}_B = \mathbf{I}
\end{equation}
which must be enforced (e.g.\ via Gram--Schmidt or quaternion
renormalisation) at every mutation site.

\paragraph{Physical role.}
$\mathbf{Q}_B$ is the \emph{pose operator} for the descriptor.
The coefficients $\{c_{\ell m}^{(k)}\}$ are defined in the bead's local
frame; $\mathbf{Q}_B$ tells the interaction engine how to rotate those
coefficients into the frame of an interaction partner.  Without
$\mathbf{Q}_B$, the descriptor has no spatial meaning.

Concretely, $\mathbf{Q}_B$ enters:
\begin{itemize}
    \item the relative rotation $R = \mathbf{Q}_A^T \mathbf{Q}_B$ used to
          transform $B$'s coefficients into $A$'s local frame
    \item torque evaluation:
          $\boldsymbol{\tau}_B = -\partial U / \partial \Omega_B$
    \item rotational equations of motion (Euler or quaternion propagation)
    \item the canonical orientation established during atomistic$\to$CG mapping
\end{itemize}

\paragraph{Update frequency.}
$\mathbf{Q}_B$ is updated \emph{every dynamics step}, in parallel with
$\mathbf{R}_B$.  After coarse-graining, the bead is an independent rigid
anisotropic object; its orientation evolves dynamically and must not be
recomputed from atomistic positions during a running simulation.

\paragraph{Ownership.}
$\mathbf{Q}_B$ belongs to the \emph{kinematic shell}, alongside
$\mathbf{R}_B$.  Together they define the bead's pose---where it is and how
it is oriented.

\paragraph{Degeneracy considerations.}
For highly symmetric or near-linear fragments the inertia tensor may have
degenerate eigenvalues, making the principal-axis assignment ambiguous.
The implementation must detect this case (eigenvalue ratio
$\lambda_i / \lambda_j < \epsilon_{\mathrm{degen}}$) and fall back to a
convention-based axis assignment to ensure reproducibility.

\subsection{$\{c_{\ell m}^{(k)}\}$: Descriptor Coefficients}

The spherical harmonic coefficients are the bead's \emph{intrinsic
identity}.  Everything else in the state is either motion or bookkeeping;
the coefficients encode what the bead actually is.

\paragraph{Physical role.}
For each channel $k \in \{\text{steric},\, \text{electrostatic},\,
\text{dispersion}\}$ the surface field is expanded as
\begin{equation}
    S_B^{(k)}(\theta,\phi) = \sum_{\ell=0}^{\ell_{\max}}
        \sum_{m=-\ell}^{\ell} c_{\ell m}^{(k)}\, Y_{\ell m}(\theta,\phi)
\end{equation}
The coefficients carry physically distinct information:
\begin{itemize}
    \item \textbf{Steric:} geometric accessibility and excluded-volume pattern
    \item \textbf{Electrostatic:} anisotropic charge-field structure
    \item \textbf{Dispersion:} anisotropic attractive strength distribution
\end{itemize}

\paragraph{Internal structure.}
The coefficient set is not a flat list.  It has a channel-indexed,
band-structured hierarchy:
\begin{equation}
    \{c_{\ell m}^{(k)}\} \;\Longrightarrow\;
    \underbrace{k}_{\text{channel}}
    \;\to\;
    \underbrace{\ell}_{\text{angular order}}
    \;\to\;
    \underbrace{m}_{\text{azimuthal mode}}
\end{equation}
This nesting is not cosmetic.  Adaptive truncation, channel activation, and
per-$\ell$ kernel evaluation all operate naturally on this structure.

\paragraph{Update frequency.}
For rigid fragments the coefficients are \emph{generated once} during
atomistic$\to$CG mapping and remain fixed thereafter.  They change only when:
\begin{itemize}
    \item the adaptive truncation loop promotes $\ell_{\max}$
    \item the underlying fragment changes (bond breaking, conformational rearrangement)
    \item the system is remapped
\end{itemize}
In a standard rigid-bead simulation, the coefficients are \emph{static
intrinsic data}, reused many times per dynamics step during pairwise
interaction evaluation.

\paragraph{Ownership.}
$\{c_{\ell m}^{(k)}\}$ belongs to the \emph{identity layer}---the layer
that defines what the bead intrinsically is, independent of where it is or
how it is currently oriented.

\subsection{$M_B$: Total Mass}

\begin{equation}
    M_B = \sum_{i \in S_B} m_i
\end{equation}

\paragraph{Physical role.}
$M_B$ is the translational inertia scale.  It enters:
\begin{itemize}
    \item Newton's second law: $M_B \ddot{\mathbf{R}}_B = \mathbf{F}_B$
    \item translational kinetic energy:
          $T_{\mathrm{trans}} = \tfrac{1}{2} M_B \dot{\mathbf{R}}_B^2$
    \item momentum conservation: $\mathbf{p}_B = M_B \dot{\mathbf{R}}_B$
    \item thermostat coupling (Langevin friction, Nos\'e--Hoover mass ratios)
\end{itemize}

\paragraph{Update frequency.}
$M_B$ is a \emph{conserved scalar} that changes only if the fragment
definition changes.  It must never drift during dynamics.

\paragraph{Ownership.}
$M_B$ belongs to the \emph{identity layer}, alongside the coefficients.
It is not part of the descriptor machinery---it describes how the bead
responds to force, not how it interacts anisotropically.

\paragraph{Separation principle.}
Descriptor coefficients tell you \emph{how the bead interacts
anisotropically}.  Mass tells you \emph{how the bead responds dynamically
to force}.  These are distinct responsibilities and must not be conflated.

\subsection{$Q_B^{\mathrm{tot}}$: Total Charge}

\begin{equation}
    Q_B^{\mathrm{tot}} = \sum_{i \in S_B} q_i
\end{equation}

\paragraph{Physical role.}
$Q_B^{\mathrm{tot}}$ is the net scalar charge of the bead.  It enters:
\begin{itemize}
    \item global charge-balance verification
    \item monopole contributions to long-range electrostatic treatments
          (Ewald summation, reaction-field corrections)
    \item system-level diagnostics and conservation checks
\end{itemize}

\paragraph{Relationship to the electrostatic channel.}
$Q_B^{\mathrm{tot}}$ and the electrostatic descriptor coefficients
$\{c_{\ell m}^{(\mathrm{elec})}\}$ are related but not identical:
\begin{itemize}
    \item $Q_B^{\mathrm{tot}}$ is the \emph{net charge}---a single scalar
    \item $c_{00}^{(\mathrm{elec})}$ is proportional to the monopole moment
          of the electrostatic surface field
    \item higher-order coefficients $c_{\ell m}^{(\mathrm{elec})}$ with
          $\ell \geq 1$ encode the \emph{angular structure} of the
          electrostatic field
\end{itemize}
Both must be retained as explicit state variables.  The scalar charge serves
conservation and long-range bookkeeping; the coefficients serve anisotropic
short-range interaction evaluation.  The channel coefficients must not
``replace'' the total charge as an explicit bead property.

\paragraph{Update frequency.}
$Q_B^{\mathrm{tot}}$ is a \emph{conserved scalar} that changes only if
the fragment is remapped.

\paragraph{Ownership.}
$Q_B^{\mathrm{tot}}$ belongs to the \emph{identity layer}.

\subsection{State Layer Architecture}

The preceding analysis reveals a natural four-layer decomposition of the
bead state.

\subsubsection{Layer A: Identity}

The intrinsic properties of the bead, generated at mapping time and
thereafter fixed unless remapping occurs:
\begin{equation}
    \mathcal{I}_B = \bigl\{ \{c_{\ell m}^{(k)}\},\; M_B,\; Q_B^{\mathrm{tot}} \bigr\}
\end{equation}

\subsubsection{Layer B: Pose}

The current spatial configuration of the bead, evolving under the equations
of motion at every dynamics step:
\begin{equation}
    \mathcal{P}_B = \bigl\{ \mathbf{R}_B,\; \mathbf{Q}_B \bigr\}
\end{equation}

\subsubsection{Layer C: Resolution}

Metadata governing the trustworthiness and complexity of the descriptor.
Not part of the mathematical state tuple, but logically attached to the
coefficient block:
\begin{itemize}
    \item active channel set: which of the three channels carry nonzero
          coefficients
    \item current $\ell_{\max}$ per channel
    \item reconstruction residual $\varepsilon_B^{(k)}$ per channel
    \item frame construction method and well-definedness flag
\end{itemize}
These determine whether the bead's descriptor is trustworthy at its current
resolution level.

\subsubsection{Layer D: Provenance}

Debugging and traceability metadata, also not part of the mathematical
state, but essential for an anti-black-box architecture:
\begin{itemize}
    \item source fragment identifier (which atomistic fragment generated
          this bead)
    \item frame method used (inertia, bond-based, fallback)
    \item whether the frame was well-defined or degenerate
    \item the residual achieved at the time of descriptor generation
    \item original atom count and bond count
\end{itemize}
Provenance does not change the physics, but it determines whether the
architecture is debuggable.

\subsubsection{Layer summary}

\begin{center}
\begin{tabular}{llll}
    \textbf{Layer} & \textbf{Contents} & \textbf{Update frequency} & \textbf{Role} \\
    \hline
    A (Identity)    & $\{c_{\ell m}^{(k)}\}$, $M_B$, $Q_B^{\mathrm{tot}}$
                    & Fixed unless remapped & What the bead is \\
    B (Pose)        & $\mathbf{R}_B$, $\mathbf{Q}_B$
                    & Every dynamics step & Where and how it is oriented \\
    C (Resolution)  & $\ell_{\max}$, active channels, residuals
                    & At adaptation events & Whether the descriptor is sufficient \\
    D (Provenance)  & Fragment ID, frame method, diagnostics
                    & At mapping time only & Traceability and debugging \\
\end{tabular}
\end{center}

\subsection{Update Frequency Classification}

The four-layer decomposition induces a strict update-frequency hierarchy
that prevents unnecessary recomputation.

\paragraph{Frequently updated (every step).}
\begin{equation}
    \mathbf{R}_B,\; \mathbf{Q}_B
\end{equation}
These are the evolving kinematic variables.  They change under force and
torque at every integration step.

\paragraph{Infrequently updated (at adaptation events).}
\begin{equation}
    \{c_{\ell m}^{(k)}\},\; \ell_{\max},\; \varepsilon_B^{(k)}
\end{equation}
For rigid fragments the coefficients remain fixed during normal dynamics.
They change only when the adaptive truncation loop promotes $\ell_{\max}$
or the underlying fragment is modified.

\paragraph{Essentially fixed (conserved unless remapped).}
\begin{equation}
    M_B,\; Q_B^{\mathrm{tot}}
\end{equation}
These must not drift during dynamics.  A change in either implies that
the underlying atomistic fragment has changed.

This classification prevents the antipattern of recomputing every piece of
bead state at every step---a design characteristic of systems built by
implementations that are afraid of structure.

\subsection{Interaction Engine View}

When evaluating the pairwise energy $U_{AB}$, the interaction engine should
receive a \emph{projected physical view} of each bead, not the entire bead
object indiscriminately.  The engine consumes:
\begin{itemize}
    \item \textbf{Relative position} from $\mathbf{R}_A$, $\mathbf{R}_B$:
          \begin{equation}
              \mathbf{r}_{AB} = \mathbf{R}_B - \mathbf{R}_A, \quad
              r = |\mathbf{r}_{AB}|
          \end{equation}
    \item \textbf{Orientation relation} from $\mathbf{Q}_A$, $\mathbf{Q}_B$:
          \begin{equation}
              R = \mathbf{Q}_A^T \mathbf{Q}_B
          \end{equation}
    \item \textbf{Descriptor content} from $\{c_{\ell m}^{(k)}\}_A$,
          $\{c_{\ell m}^{(k)}\}_B$
\end{itemize}

The engine does \emph{not} need $M_B$ or $Q_B^{\mathrm{tot}}$ directly,
except where a separate monopole term or long-range electrostatic treatment
requires the scalar charge.

This separation ensures that the interaction engine operates on the
physically relevant projection of the bead state, not on an undifferentiated
data dump.

\subsection{Dynamics Engine View}

The motion engine consumes a different projection of the bead state:
\begin{itemize}
    \item $\mathbf{R}_B$: position to be integrated
    \item $\mathbf{Q}_B$: orientation to be integrated
    \item $M_B$: translational inertia for force$\to$acceleration conversion
    \item $\mathbf{I}_B$: rotational inertia for torque$\to$angular acceleration
          conversion (see below)
    \item $\mathbf{F}_B$, $\boldsymbol{\tau}_B$: force and torque from the
          interaction engine
\end{itemize}

The dynamics engine does \emph{not} need to know about descriptor generation
details, channel sampling procedures, or spherical harmonic projection
methodology.  It needs the resulting force and torque, not a reconstruction
of how the steric channel was projected onto harmonics.

This separation is healthy: the interaction engine converts descriptor
content into forces and torques; the dynamics engine converts forces and
torques into motion.

\subsection{Rotational Inertia Extension}

The formal state as written is sufficient for descriptor representation and
interaction evaluation.  However, for complete rigid-body rotational
dynamics the bead requires one additional quantity: the \emph{rotational
inertia tensor}.

For a rigid bead composed of atoms $i \in S_B$ with positions
$\boldsymbol{\delta}_i = \mathbf{r}_i - \mathbf{R}_B$ relative to the
centre of mass:
\begin{equation}
    \mathbf{I}_B = \sum_{i \in S_B} m_i
    \left( |\boldsymbol{\delta}_i|^2 \,\mathbf{I}_{3}
           - \boldsymbol{\delta}_i \boldsymbol{\delta}_i^T \right)
\end{equation}

Diagonalisation in the principal-axis frame yields:
\begin{equation}
    \mathbf{I}_B = \mathbf{Q}_B \,
    \mathrm{diag}(I_1, I_2, I_3) \,
    \mathbf{Q}_B^T
\end{equation}

The principal moments $(I_1, I_2, I_3)$ are scalars that, like $M_B$,
are generated at mapping time and remain fixed unless the fragment changes.

With this extension, Euler's rotational equations of motion become:
\begin{equation}
    I_\alpha \dot{\omega}_\alpha
    = \tau_\alpha^{\mathrm{body}}
    - \epsilon_{\alpha\beta\gamma}\, \omega_\beta\, I_\gamma\, \omega_\gamma
    \qquad (\alpha = 1,2,3)
\end{equation}
where $\omega_\alpha$ are the body-frame angular velocity components and
$\tau_\alpha^{\mathrm{body}}$ are the body-frame torque components.

The theoretically complete dynamic bead state is therefore:
\begin{equation}
\boxed{
    B = \bigl\{
        \mathbf{R}_B,\;
        \mathbf{Q}_B,\;
        \{c_{\ell m}^{(k)}\},\;
        M_B,\;
        Q_B^{\mathrm{tot}},\;
        \mathbf{I}_B
    \bigr\}
}
\end{equation}

This is the state of a genuinely physical anisotropic bead---not merely a
clever interaction shell, but an object with full translational and
rotational dynamics grounded in the inertia of its underlying atomistic
fragment.

\subsection{Conceptual Summary}

The bead state admits a clean three-way separation:
\begin{equation}
    \underbrace{\mathbf{R}_B,\; \mathbf{Q}_B}_{\text{where and how it is oriented}}
    \quad\Big|\quad
    \underbrace{\{c_{\ell m}^{(k)}\}}_{\text{what it intrinsically is}}
    \quad\Big|\quad
    \underbrace{M_B,\; Q_B^{\mathrm{tot}},\; \mathbf{I}_B}_{\text{how it responds to force and torque}}
\end{equation}

\noindent
This separation must not be blurred:
\begin{itemize}
    \item The \textbf{pose} ($\mathbf{R}_B$, $\mathbf{Q}_B$) is dynamic---it
          evolves every step.
    \item The \textbf{descriptor} ($\{c_{\ell m}^{(k)}\}$) is intrinsic---it
          encodes the bead's angular identity and changes only under
          remapping or adaptation.
    \item The \textbf{inertial properties} ($M_B$, $Q_B^{\mathrm{tot}}$,
          $\mathbf{I}_B$) are conserved scalars inherited from the atomistic
          fragment.
\end{itemize}

An effective implementation treats these three groups with distinct storage
strategies, distinct access patterns, and distinct update schedules.  The
interaction engine sees pose and descriptor.  The dynamics engine sees pose
and inertia.  Neither needs to see the other's internal details.  That is
the right architecture.

\section{Multi-Layer Bead Visualization Model}

The canonical bead state provides a complete mathematical description of the
coarse-grained object: position, orientation, descriptor coefficients, mass,
charge, and inertia.  However, for research-oriented inspection, debugging,
and scientific reporting, the bead must also be \emph{visually decomposable}
into distinct representational layers.  This section formalises the
multi-layer bead visualisation model, in which each layer reveals a different
aspect of the bead's inherited internal structure.

\subsection{Layer Architecture Overview}

The bead's visual representation is decomposed into three layers, each
serving a distinct interpretive purpose:

\begin{center}
\begin{tabular}{clp{7cm}}
    \textbf{Layer} & \textbf{Name} & \textbf{Purpose} \\
    \hline
    1 & Effective Shell & Outer boundary sphere encoding effective radius,
        mapping confidence, and structural classification (aromatic, cyclic,
        organometallic) \\
    2 & Surface-State Field & Directional effective state inherited from the
        decomposed molecular fragment, rendered as a spherical heatmap \\
    3 & Internal Structural Reference & Original atomistic geometry preserved
        inside the bead volume for interpretability and provenance \\
\end{tabular}
\end{center}

\noindent
Layers are independent and composable: any subset may be rendered
simultaneously.  The default inspection sequence descends from Layer~1
(overview) through Layer~2 (anisotropic identity) to Layer~3 (atomistic
provenance).

\subsection{Layer 1: Effective Shell (review)}

Layer~1 is the bead's outer boundary, already defined in the preceding
sections.  It represents the effective extent of the bead in space and
encodes:
\begin{itemize}
    \item effective radius (from probe sampling or mass estimate)
    \item structural classification (aromatic, cyclic, organometallic)
    \item mapping confidence (opacity modulation)
    \item frame validity (axes overlaid when the local frame is well-defined)
\end{itemize}

Layer~1 is the default scene-level representation.  It carries no
directional field information.

\subsection{Layer 2: Effective Surface-State Layer}

Layer~2 is the bead's directional surface-state representation, defined on
a spherical envelope centred at the bead reference point and oriented by the
bead-local frame $\mathbf{Q}_B$.

\subsubsection{Physical meaning}

This layer does not represent literal atomic boundary geometry and does not
primarily represent incoming external forces.  Instead, it is the
\emph{effective spherical projection} of the underlying fragment's
internally inherited structure, after decomposition into the bead state.

It encodes the directional residue of the original molecular component,
including effects inherited from:
\begin{itemize}
    \item mass distribution
    \item bonded geometry
    \item local charge arrangement
    \item electron-density asymmetry
    \item anisotropic stiffness, accessibility, or cohesive tendency
    \item other channel-specific reduced descriptors
\end{itemize}

\subsubsection{Formal definition}

Layer~2 is a scalar or multi-channel field over the sphere,
\begin{equation}
    S^{(k)}(\theta, \phi)
\end{equation}
where each channel $k$ represents one reduced internal state quantity
(steric, electrostatic, dispersive, or a synthesised effective-force
measure).  The field is evaluated from the spherical harmonic expansion:
\begin{equation}
    S^{(k)}(\theta, \phi)
    = \sum_{\ell=0}^{\ell_{\max}}
      \sum_{m=-\ell}^{\ell}
      c_{\ell m}^{(k)}\, Y_{\ell m}(\theta, \phi)
\end{equation}

The purpose of Layer~2 is to make visible the statement:

\begin{quote}
``What internal directional structure survives after the molecular fragment
has been compressed into a bead?''
\end{quote}

\subsubsection{Functional role}

Layer~2 serves as the bead's primary visible anisotropic identity.  It is
the intermediate object between:
\begin{itemize}
    \item the hidden source fragment (Layer~3)
    \item the bead's later external interaction behaviour
\end{itemize}

The causal chain is:
\begin{equation}
    \text{underlying molecular fragment}
    \;\xrightarrow{\text{reduction}}\;
    \text{directional state}
    \;\xrightarrow{\text{projection}}\;
    \text{surface-state field}
    \;\xrightarrow{\text{coupling}}\;
    \text{effective bead interactions}
\end{equation}

\subsubsection{Visual form}

In rendering, Layer~2 appears as:
\begin{itemize}
    \item a spherical shell centred at the bead origin
    \item sampled over angular coordinates $(\theta, \phi)$ on a regular grid
    \item coloured by field magnitude or signed value using a continuous
          colourmap
    \item optionally shown with transparency
    \item optionally shown with radial exaggeration for emphasis, though
          spherical geometry should remain the default
\end{itemize}

This shell is therefore a \emph{field visualisation}, not a hard-body
surface.  The rendering uses quad-strip tessellation over the
$(\theta, \phi)$ grid, with per-vertex colour interpolation derived from
the active descriptor channel.

\subsubsection{Channel selection}

The user may select which channel to display:
\begin{itemize}
    \item $k = 0$: Steric (geometric accessibility, excluded volume)
    \item $k = 1$: Electrostatic (charge-field structure)
    \item $k = 2$: Dispersion (attractive strength distribution)
    \item Synthesised: a weighted combination
          $S_{\mathrm{eff}} = w_s S^{(0)} + w_e S^{(1)} + w_d S^{(2)}$
\end{itemize}

\subsection{Layer 3: Internal Structural Reference Layer}

Layer~3 is the bead's internal structural reference representation.  It
preserves a view of the original molecular component, or a simplified
retained substructure, inside the bead volume.

\subsubsection{Physical meaning}

This layer is not the active coarse-grained state itself.  It is a
\emph{reference and interpretability layer} used to show the structural
source from which Layer~2 was derived.

It may include:
\begin{itemize}
    \item atom positions (from the source \texttt{atomistic::State})
    \item bond connectivity (from the \texttt{FragmentView} bond records)
    \item ring or fragment geometry
    \item local orientation markers
    \item optional partial charge or element colouring
\end{itemize}

\subsubsection{Functional role}

Layer~3 answers the question:
\begin{quote}
``What original internal structure gave rise to the surface pattern
shown on Layer~2?''
\end{quote}

This layer provides traceability from the coarse-grained representation
back to the atomistic source-of-truth, in accordance with the
anti-black-box design philosophy.

\subsubsection{Connection to the atomistic namespace}

Layer~3 uses the project's native atomistic representation:
\begin{itemize}
    \item \texttt{atomistic::Vec3} for atom positions
    \item \texttt{atomistic::AtomRecord} for element identity, charge, and
          classification flags
    \item \texttt{atomistic::BondRecord} for connectivity and bond order
    \item \texttt{atomistic::LocalFrame} for the fragment's structural
          orientation
\end{itemize}

The visual record stores a \texttt{SourceFragmentMiniature} that packages
these data into a rendering-ready format without requiring the renderer to
import or re-traverse raw atomistic state.

\subsubsection{Visual form}

In rendering, Layer~3 appears as:
\begin{itemize}
    \item small spheres at each atom position, coloured by element (CPK
          convention)
    \item stick bonds between bonded atoms, with line width modulated by
          bond order
    \item optional orientation markers (local frame axes of the fragment)
    \item the entire miniature is positioned inside the Layer~1 shell at
          the fragment's centre of mass
\end{itemize}

\subsection{Layer Composition}

The three layers compose into a set of view modes:

\begin{center}
\begin{tabular}{lcccp{4.5cm}}
    \textbf{View Mode} & \textbf{L1} & \textbf{L2} & \textbf{L3}
    & \textbf{Description} \\
    \hline
    Shell           & \checkmark &            &            & Overview only \\
    Scaffold        & ghost      &            &            & Shell + frame + solved
                                                            directions \\
    SurfaceState    & ghost      & \checkmark &            & Spherical heatmap of
                                                            effective state \\
    InternalRef     & ghost      &            & \checkmark & Atomistic fragment inside
                                                            ghost shell \\
    Skeleton        & ghost      &            & \checkmark & Alias for InternalRef \\
    Cutaway         & clipped    & clipped    &            & Hemisphere section \\
    Residual        & ghost      &            &            & Error bands and halos \\
    Comparison      & ghost      & \checkmark & \checkmark & Layers 2 and 3
                                                            overlaid \\
\end{tabular}
\end{center}

\noindent
``ghost'' denotes a semi-transparent shell.  ``clipped'' denotes a shell
intersected by a user-controlled clipping plane.

\subsection{Anti-Black-Box Properties}

The multi-layer model satisfies the project's anti-black-box requirements:
\begin{itemize}
    \item Layer~1 makes the bead's spatial extent and classification
          explicitly visible.
    \item Layer~2 makes the bead's anisotropic identity inspectable at every
          angular direction.
    \item Layer~3 makes the atomistic provenance of every descriptor
          coefficient traceable to specific atoms and bonds.
    \item All three layers are independently toggleable, composable, and
          deterministic: the same bead state always produces the same visual
          output.
\end{itemize}

\section{Conclusion}

This document has developed the anisotropic surface-mapped bead from initial
concept to a formally specified, architecturally decomposed physical object.

The progression is intentional and cumulative:
\begin{enumerate}
    \item \textbf{Atomistic-to-CG mapping} established the centre-of-mass,
          local frame, and aggregate scalars that anchor every bead to its
          underlying fragment.
    \item \textbf{Surface descriptor representation} replaced isotropic
          point particles with directional surface fields, expanded in
          spherical harmonics for compact storage and efficient evaluation.
    \item \textbf{Orientation-coupled interactions} introduced relative
          alignment dependence, torque, and rotational dynamics---turning
          the bead from a passive site into a dynamical object with angular
          degrees of freedom.
    \item \textbf{Multi-channel descriptor enrichment} decomposed the
          surface field into steric, electrostatic, and dispersion channels,
          enabling physically distinct interaction mechanisms to be resolved
          independently.
    \item \textbf{Unified descriptor strategy} eliminated discrete model
          switching in favour of a single continuous formalism with
          residual-driven adaptive truncation.
    \item \textbf{Atomistic preparation layer} formalised the scale boundary
          via the fragment view, ensuring that every descriptor coefficient
          is traceable to specific atoms and that invalid inputs are rejected
          with explicit status codes.
    \item \textbf{Implementation specification} fixed the per-$\ell$
          channel kernels, the Wigner-$D$ rotation pipeline, and the
          adaptive refinement loop into a rigid, deterministic architecture.
    \item \textbf{Formal bead state decomposition} analysed each variable in
          the canonical tuple, established the four-layer architecture
          (identity, pose, resolution, provenance), classified update
          frequencies, defined the interaction and dynamics engine views,
          and extended the state to include rotational inertia for complete
          rigid-body dynamics.
\end{enumerate}

The result is a bead defined by the complete state
\begin{equation}
    B = \bigl\{
        \mathbf{R}_B,\;
        \mathbf{Q}_B,\;
        \{c_{\ell m}^{(k)}\},\;
        M_B,\;
        Q_B^{\mathrm{tot}},\;
        \mathbf{I}_B
    \bigr\}
\end{equation}
with clear ownership, update rules, and numerical roles for every term.

\subsection{Scope Boundary}

This document deliberately stops at the \emph{isolated bead}---its internal
structure, its descriptor, its interactions with other beads, and its
rigid-body dynamics.  The bead as defined here is a self-contained physical
object: it knows what it is, how it is oriented, how it interacts, and how
it moves.

What the bead does \emph{not} yet know is its environment.  It has no
awareness of local density, coordination, orientational order, or collective
state.  It cannot distinguish between being in a gas, a liquid, a crystal,
or at an interface.  It cannot respond to crowding, stiffen under
compression, or relax into alignment with its neighbours.

These capabilities require a second layer of state---environment-responsive
observables and slow internal variables---that couples the bead to its local
surroundings.  This extension is the subject of a companion document:

\begin{quote}
    \emph{Environment-Responsive Coarse-Grained Bead Dynamics} \\
    (VSEPR-SIM Project, \texttt{section\_environment\_responsive\_beads.tex})
\end{quote}

\noindent
That document introduces the extended bead state
$\mathbb{B}_B = \{B,\, X_B\}$, where $X_B$ contains environment
observables (local density, coordination, orientational order) and a slow
internal state variable $\eta_B$ with memory and relaxation dynamics.  The
coupling of $X_B$ back into the interaction kernels enables
density-dependent behaviour, spontaneous ordering, hysteresis, and
metastability---the transition from ``the model works'' to ``the system
behaves.''

\end{document}
